
\subsubsection*{4.2.2}
\phantomsection
\label{sga3:VIIB:4.2.2}

Suppose that \(\mathbf L\) is a flat \(p\)-Lie algebra over
\(\mathbf O_k\). Taking VII\textsubscript{A}, \SGARefLink{sga3:VIIA:5.3.3}{5.3.3} into account,
the argument of part \textup{\SGARefLink{sga3:vii:paragraph:5-1:equation:i}{(i)}} of \SGARefLink{sga3:VIIB:2.6.2:prop}{Proposition~2.6.2} shows that
\(U_p(\mathbf L)\) is flat over \(\mathbf O_k\). Consequently,
\(\mathscr G_p(\mathbf L)\) is a formal \(k\)-group topologically
flat over \(k\), with \(U_p(\mathbf L)\) as its covariant
\(\mathbf O_k\)-bialgebra.

\footnote{\label{sga3:VIIB:nde:135}The editors supplied the
details that follow.}
By the proof of \SGARefLink{sga3:VIIB:2.6.2:prop}{2.6.2}\textup{\SGARefLink{sga3:VIIB:2.6.2:prop:item:iii}{(iii)}}, for every finite-length
\(k\)-algebra \(C\), the Lie algebra
\(\operatorname{Lie}(\mathscr G_p(\mathbf L))(C)\) is the set of
primitive elements of
\[
  U_p(\mathbf L)(C)
  =
  U_p\bigl(\mathbf L(C)\bigr);
\]
see also VII\textsubscript{A}, \SGARefLink{sga3:VIIA:3.2.3}{3.2.3}. Moreover, by
VII\textsubscript{A}, \SGARefLink{sga3:VIIA:5.5.3}{5.5.3}, the canonical morphism
\(\mathbf L\to U_p(\mathbf L)\) induces an isomorphism of
\(p\)-Lie algebras
\[
  \tau_{p,\mathbf L}:
  \mathbf L
  \xrightarrow{\ \sim\ }
  \operatorname{Lie}\bigl(\mathscr G_p(\mathbf L)\bigr).
\]
Compare this with \SGARefLink{sga3:VIIB:3.1:lemma}{3.1} in characteristic \(0\).

The formal group \(\mathscr G_p(\mathbf L)\) can be characterized
by a universal property. Every morphism
\[
  h:\mathscr G_p(\mathbf L)\longrightarrow G
\]
to a formal group \(G\) induces a morphism
\[
  \operatorname{Lie}(h):
  \operatorname{Lie}\bigl(\mathscr G_p(\mathbf L)\bigr)
  \longrightarrow
  \operatorname{Lie}(G).
\]
Composing this morphism with \(\tau_{p,\mathbf L}\) gives
\[
  h':\mathbf L\longrightarrow\operatorname{Lie}(G).
\]

\medskip
\noindent\textbf{Proposition.}
\phantomsection
\label{sga3:VIIB:4.2.2:universal-property}
{\itshape Let \(k\) be a pseudocompact ring of characteristic
\(p>0\), let \(G\) be a formal \(k\)-group, and let
\(\mathbf L\) be a flat \(p\)-Lie algebra over \(\mathbf O_k\).
The map \(h\mapsto h'\) defined above is a bijection
\[
  \operatorname{Hom}_{\operatorname{Grf}/k}
    \bigl(\mathscr G_p(\mathbf L),G\bigr)
  \xrightarrow{\ \sim\ }
  \operatorname{Hom}_{p\text{-}\operatorname{Lie}}
    \bigl(\mathbf L,\operatorname{Lie}(G)\bigr),
\]
where the right-hand side is the set of morphisms of \(p\)-Lie
algebras from \(\mathbf L\) to \(\operatorname{Lie}(G)\).}

\textit{Proof.}
We may immediately reduce to the case in which \(k\) is Artinian.
Put
\[
  L=\mathbf L(k).
\]
By \SGARefLink{sga3:VIIB:2.3.1}{2.3.1},
\[
  \operatorname{Hom}_{\operatorname{Grf}/k}
    \bigl(\mathscr G_p(\mathbf L),G\bigr)
\]
is in bijection with the set of unital algebra morphisms
\[
  h:U_p(L)\longrightarrow\mathbf H(G)
\]
for which the following diagram commutes:
\begin{center}
  \phantomsection
  \label{sga3:VIIB:4.2.2:restricted-universal-diagram}
  \begin{tikzcd}[row sep=1.2em,column sep=3.9em]
U_p(L) \arrow[r,"h"] \arrow[dd,"\Delta_{U_p(L)}"']
  & \mathbf H(G) \arrow[dr,"\delta_G"] & {} \\
{} & {} & \mathbf H(G\times G) \\
U_p(L)\otimes U_p(L) \arrow[r,"h\otimes h"]
  & \mathbf H(G)\otimes\mathbf H(G)
      \arrow[ur,"\varphi_G"']
  & {.}
\end{tikzcd}

\end{center}
The morphism \(h\) is determined by its restriction to \(L\),
which is a morphism of \(p\)-Lie algebras from \(L\) to the
underlying \(p\)-Lie algebra of \(\mathbf H(G)\). The commutativity
of the diagram says precisely that \(h\) sends \(L\) into the
subset of \(\mathbf H(G)\) consisting of those \(x\) such that
\[
  \delta_G(x)
  =
  \varphi_G(x\otimes1+1\otimes x).
\]
This subset is \(\operatorname{Lie}(G)\), by \SGARefLink{sga3:VIIB:2.6}{2.6}\textup{(c)}.
The proposition follows.
